\documentclass[tikz,border=10pt]{standalone}
\usepackage{amsmath}
\usepackage{amssymb}
\usepackage{tikz}
\usetikzlibrary{shapes.geometric, arrows.meta, positioning, fit, backgrounds, shadows, calc}

\begin{document}

\tikzset{
    box/.style={
        rectangle,
        rounded corners=3mm,
        draw=#1!70,
        fill=#1!15,
        thick,
        minimum height=1cm,
        text width=2.5cm,
        align=center,
        font=\small\bfseries,
        drop shadow
    },
    data/.style={
        cylinder,
        draw=blue!60,
        fill=blue!10,
        minimum height=1.2cm,
        minimum width=1.5cm,
        shape border rotate=90,
        font=\tiny\bfseries,
        drop shadow
    },
    arrow/.style={
        -Stealth,
        very thick,
        draw=gray!70
    }
}

\begin{tikzpicture}[node distance=1.2cm and 1cm]

    % Title
    \node[font=\Large\bfseries, text=blue!80!black] (title) {Instance-Based Domain Adaptation};

    % Source and target
    \node[data, below left=1cm and 2cm of title] (source) {Source\\Data\\(Labeled)};
    \node[data, below right=1cm and 2cm of title] (target) {Target\\Data\\(Unlabeled)};

    % Pre-trained model
    \node[box=blue, below=1.5cm of title] (pretrained) {Pre-trained\\Model};

    % Instance weighting
    \node[box=orange, below left=of pretrained] (weighting) {Instance\\Weighting};

    % Self-training
    \node[box=purple, below right=of pretrained] (selftrain) {Self-Training\\(Pseudo-labels)};

    % Adapted model
    \node[box=green, below=2cm of pretrained] (adapted) {Adapted\\Model};

    % Test time adaptation
    \node[box=teal, below=of adapted] (tta) {Test-Time\\Adaptation};

    % Final predictions
    \node[box=green!70!black, below=of tta] (output) {Predictions};

    % Arrows
    \draw[arrow] (source) -- (pretrained);
    \draw[arrow] (pretrained) -- (weighting);
    \draw[arrow] (pretrained) -- (selftrain);
    \draw[arrow] (target) -- (selftrain);
    \draw[arrow] (weighting) -- (adapted);
    \draw[arrow] (selftrain) -- (adapted);
    \draw[arrow] (adapted) -- (tta);
    \draw[arrow] (target) -- (tta);
    \draw[arrow] (tta) -- (output);

    % Annotations
    \node[font=\scriptsize, right=0.1cm of weighting, text width=2.2cm, align=left] {
        Reweight\\samples by\\importance
    };

    \node[font=\scriptsize, left=0.1cm of selftrain, text width=2.5cm, align=right] {
        Generate labels\\for target data
    };

    % Benefits
    \node[box=green, below=0.6cm of output, minimum width=8cm, text width=7.5cm, font=\scriptsize\bfseries] (benefits) {
        \textbf{Advantages:} No target labels required, adapts at test time, flexible\\
        \textbf{Limitations:} Requires confidence calibration, sensitive to distribution shift
    };

\end{tikzpicture}

\end{document}
